% Source: https://artofproblemsolving.com/community/c3h2489013_problem_of_divisor
% Stable identity: aops-topic-2489013-q01
Note $xy-1\ge1$ is needed (for $x=y=1$, $xy-1=0$ divides only $0$). Any common divisor of $xy-1$ and $x$ divides $xy-(xy-1)=1$, so $\gcd(xy-1,x)=1$. Since $x^3+x=x\left(x^2+1\right)$, the condition is equivalent to
\[xy-1\mid x^2+1.\]
Similarly $\gcd(xy-1,y)=1$, and $y^2\left(x^2+1\right)=(xy)^2+y^2\equiv1+y^2\pmod{xy-1}$, so the condition is also equivalent to $xy-1\mid y^2+1$.

\textbf{Bounding.} A positive divisor is at most the number, so $xy-1\le x^2+1$ and $xy-1\le y^2+1$, i.e.
\[y\le x+\dfrac2x,\qquad x\le y+\dfrac2y.\]
If $x\ge3$ and $y\ge3$, these give $y<x+1$ and $x<y+1$, so $x=y$; then $x^2-1\mid x^2+1$, so $x^2-1\mid2$, impossible for $x\ge3$. Hence $x\le2$ or $y\le2$.
\begin{itemize}
\item $x=1$: $y-1\mid2$, so $y=2$ or $y=3$.
\item $x=2$: $2y-1\mid5$, so $y=1$ or $y=3$.
\item $y=1$, $x\ge3$: $x-1\mid x^2+1=(x-1)(x+1)+2$, so $x-1\mid2$ and $x=3$.
\item $y=2$, $x\ge3$: $2x-1\mid x^2+1$, hence $2x-1\mid4x^2+4=(2x-1)(2x+1)+5$, so $2x-1\mid5$ and $x=3$.
\end{itemize}
All of these work, so the solutions are
\[\boxed{(x,y)=(1,2),\ (1,3),\ (2,1),\ (2,3),\ (3,1),\ (3,2)}.\]
